\documentclass[11pt,a4paper]{article}

\usepackage[margin=2.2cm]{geometry}
\usepackage{array}
\usepackage{longtable}
\usepackage{xcolor}
\usepackage{hyperref}
\usepackage{enumitem}

\definecolor{pmoteal}{HTML}{2F6F73}
\definecolor{pmoamber}{HTML}{D99A2B}
\definecolor{pmogray}{HTML}{405059}

\hypersetup{
  colorlinks=true,
  linkcolor=pmoteal,
  urlcolor=pmoteal
}

\setlist[itemize]{leftmargin=1.4em}
\setlist[enumerate]{leftmargin=1.6em}
\renewcommand{\arraystretch}{1.25}

\title{\textbf{Lecture Syllabus}\\From Plan to Control: Managing a Project Like a Real PMO}
\author{AI-Assisted Gamified Project Management Module}
\date{}

\begin{document}
\maketitle

\section*{Lecture Information}

\begin{tabular}{>{\bfseries}p{4cm}p{10.5cm}}
Module type & Gamified AI-assisted university lecture \\
Target audience & Third- or fourth-year undergraduate students, or Licence/Master students in engineering, computer science, business, or management \\
Duration & One 3-hour session, or two 90-minute sessions \\
Main case & Smart Campus Mobile App project \\
License & CC-BY-NC recommended \\
\end{tabular}

\section*{Lecture Description}

In this lecture, students act as a project management office team responsible for delivering a Smart Campus Mobile App in 12 weeks. The project has a limited budget, a small technical team, a required demo by week 6, and a major integration risk with the university database.

Students do not only listen to project management concepts. They define project scope, identify stakeholders, build a Work Breakdown Structure, prepare a schedule, analyze dashboard indicators, face a project crisis, and present a corrective action plan.

\section*{Learning Objectives}

By the end of the lecture, students will be able to:

\begin{enumerate}
  \item Explain the project management cycle: initiation, planning, execution, monitoring, control, and closing.
  \item Build a deliverable-oriented Work Breakdown Structure for a realistic project.
  \item Convert a WBS into tasks, dependencies, milestones, and deliverables.
  \item Estimate duration, cost, resources, and risks.
  \item Optimize project execution by adjusting resources, sequencing, and priorities.
  \item Track progress using planned value, earned value, actual cost, SV, CV, SPI, and CPI.
  \item Diagnose project problems from dashboard data.
  \item Propose corrective actions when a project is late, over budget, or technically unstable.
  \item Communicate project status in a professional PMO-style briefing.
\end{enumerate}

\section*{Case Scenario}

The university wants to build a Smart Campus Mobile App in 12 weeks.

\subsection*{App Features}

\begin{itemize}
  \item Student login
  \item Course schedule
  \item Exam calendar
  \item Notifications
  \item Room reservation
  \item Complaint/request system
  \item Admin dashboard
  \item Pilot deployment
\end{itemize}

\subsection*{Project Constraints}

\begin{tabular}{>{\bfseries}p{4cm}p{10.5cm}}
Budget & 40,000 USD \\
Duration & 12 weeks \\
Team & 1 project manager, 2 backend developers, 2 mobile developers, 1 UI/UX designer, 1 QA tester \\
Main risk & Integration delay with the university database \\
Sponsor pressure & First demo required by week 6 \\
\end{tabular}

\section*{Gamified Progression}

The lecture uses five levels. Each level contains a mission, a deliverable, and a badge.

\begin{longtable}{>{\bfseries}p{1.6cm}p{3.2cm}p{5.2cm}p{3.4cm}}
\textbf{Level} & \textbf{Name} & \textbf{Mission} & \textbf{Badge} \\
\hline
1 & Project Launch & Define scope, stakeholders, constraints, and success criteria & Scope Master \\
2 & WBS Builder & Build a deliverable-oriented WBS & WBS Architect \\
3 & Schedule Rescue & Compress a 15-week plan into a realistic 12-week schedule & Schedule Optimizer \\
4 & Dashboard Detective & Analyze week 6 dashboard data and diagnose project health & Dashboard Detective \\
5 & Crisis Controller & Respond to the week 8 crisis and present a recovery briefing & Risk Analyst / Crisis Controller \\
\end{longtable}

\section*{Lecture Schedule}

\begin{longtable}{>{\bfseries}p{2.5cm}p{5.1cm}p{4cm}p{3cm}}
\textbf{Time} & \textbf{Activity} & \textbf{Learning Type} & \textbf{Tools} \\
\hline
0--15 min & Opening scenario and decision poll & Acquisition + discussion & Genially, Mentimeter/Wooclap \\
15--45 min & Planning concepts and project charter & Acquisition + inquiry & Slides, NotebookLM \\
45--80 min & WBS workshop & Production + collaboration & Worksheet, Canva, Mermaid/draw.io \\
80--90 min & Break & Pause & -- \\
90--120 min & Schedule optimization challenge & Practice + collaboration & Google Sheets/Excel \\
120--155 min & Dashboard and KPI analysis & Practice + inquiry & Google Sheets, Quizizz \\
155--175 min & Week 8 crisis simulation & Discussion + production & Genially, group briefing \\
175--180 min & Exit ticket & Reflection & Google Forms or paper \\
\end{longtable}

\section*{AI Tools Used}

\begin{longtable}{>{\bfseries}p{4.5cm}p{10cm}}
\textbf{Tool} & \textbf{Role in the Lecture} \\
\hline
NotebookLM & Source-grounded summaries of project management concepts from uploaded course material \\
ChatGPT / Gemini / Claude & Scenario writing, prompts, crisis cards, quiz questions, and feedback \\
Learning Designer & Main module structure, activity types, export, and public link \\
Gamma and Canva & Slide draft generation and final visual design \\
Napkin AI / Mermaid / draw.io & WBS, Gantt, RACI, risk matrix, and control-loop diagrams \\
Genially & Branching simulation and interactive decision screens \\
Quizizz / QuestionWell / MagicSchool & Formative quizzes and instant feedback \\
Google Sheets / Excel & Schedule, budget, earned value, and dashboard simulation \\
\end{longtable}

\section*{Student Deliverables}

Each group submits or presents:

\begin{enumerate}
  \item Project charter summary
  \item Stakeholder list and success criteria
  \item Deliverable-oriented WBS
  \item Optimized 12-week schedule
  \item Risk register or risk notes
  \item Dashboard diagnosis
  \item Corrective action plan
  \item Three-minute PMO recovery briefing
\end{enumerate}

\section*{Assessment}

\begin{longtable}{>{\bfseries}p{6cm}p{3cm}p{5.3cm}}
\textbf{Criterion} & \textbf{Weight} & \textbf{Focus} \\
\hline
Scope clarity & 15\% & Stakeholders, constraints, and measurable success criteria \\
WBS quality & 20\% & Deliverable orientation, completeness, and logical decomposition \\
Schedule and dependencies & 20\% & Realistic sequencing, dependencies, and optimization \\
Risk analysis & 15\% & Identification of risks and relevant mitigation actions \\
Tracking and control logic & 20\% & Correct dashboard interpretation and corrective action reasoning \\
Professional presentation & 10\% & Clear, concise, and convincing PMO briefing \\
\end{longtable}

\section*{Preparation}

Before the lecture, students should review basic definitions of project scope, stakeholders, WBS, milestone, risk, and dashboard. No advanced project management software experience is required.

\section*{Exit Ticket}

At the end of the lecture, students answer:

\begin{enumerate}
  \item What is one difference between monitoring and controlling?
  \item What was the most important corrective action in your group plan?
  \item How did AI help your team during the activity?
  \item What should a project manager never hide from the sponsor?
\end{enumerate}

\end{document}
